% Source: https://artofproblemsolving.com/community/c3h3185222_2018_abmc_speed_contest__actonboxborough_math_competition
% Stable identity: aops-topic-3185222-q24
Scale the triangle so that its inradius is $1$. Set $u=\sqrt3$, $v=\sqrt2$, and $L=u+v+1$. Coordinates realizing the angles are $A=(0,0)$, $B=(L,0)$, $C=\left(\dfrac{L}{u+1},\dfrac{uL}{u+1}\right)$. The bisector of $\angle A$ has slope $\dfrac{1}{u}$; together with distance $1$ from $AB$ this gives $I=(u,1)$. The bisector at $B$ is consistent since $L-u=v+1=\cot22.5^\circ$. Since $BC$ has slope $-1$, the altitude from $A$ is $y=x$; the altitude from $C$ is vertical. Thus $H=\left(\dfrac{L}{u+1},\dfrac{L}{u+1}\right)$. Equating distances to $A,B,C$ gives $O=\left(\dfrac{L}{2},L\left(1-\dfrac{u}{2}\right)\right)$. Let $U=O-I$ and $V=H-I$. Using $u^2=3,v^2=2$ gives $U=\dfrac{1}{2}(1+v-u,u+2v-uv-3)$ and $V=\dfrac{u-1}{2}(v-2,v)$. Direct multiplication gives $U\cdot V=(u-1)(1-v)<0$ and $\det(U,V)=(u-1)(2-u)(v-1)>0$. Therefore $\theta=\angle OIH$ is obtuse and $\tan(180^\circ-\theta)=\dfrac{\det(U,V)}{-U\cdot V}=2-\sqrt3=\tan15^\circ$. Hence $\boxed{\theta=165^\circ}$.
\begin{asy}
usepackage("amsmath");
size(7cm);
real u=sqrt(3),v=sqrt(2),L=u+v+1; pair A=(0,0),B=(L,0),C=(L/(u+1),u*L/(u+1)),I=(u,1),H=(L/(u+1),L/(u+1)),O=(L/2,L*(1-u/2)); draw(A--B--C--cycle); draw(O--I--H,blue); draw(A--H^^C--(C.x,0),dashed); dot("$A$",A,SW);dot("$B$",B,SE);dot("$C$",C,N);dot("$I$",I,E);dot("$H$",H,W);dot("$O$",O,SE); label("$60^\circ$",A+(0.55,0.25));label("$45^\circ$",B+(-0.65,0.2));
\end{asy}
